\documentclass[conference]{IEEEtran}
\IEEEoverridecommandlockouts
\usepackage{cite}
\usepackage{amsmath,amssymb,amsfonts}
\usepackage{algorithmic}
\usepackage{graphicx}
\usepackage{textcomp}
\usepackage{xcolor}
\usepackage{booktabs}
\usepackage{multirow}
\def\BibTeX{{\rm B\kern-.05em{\sc i\kern-.025em b}\kern-.08em
    T\kern-.1667em\lower.7ex\hbox{E}\kern-.125emX}}
\begin{document}

\title{Enhancing 3D Brain Tumor Segmentation via Spatial Attention
Networks, Schedule-Free Optimization, Diffusion-Based Augmentation,
and Greedy Model Souping}

\author{
\IEEEauthorblockN{Omprakash Pugazhendhi}
\IEEEauthorblockA{
  \textit{Dept. of Computer Science and Engineering}\\
  \textit{Vellore Institute of Technology}\\
  Chennai, India\\
  omprakash.2021@vitstudent.ac.in}
}

\maketitle

\begin{abstract}
Accurate segmentation of brain tumors in multimodal MRI is
essential for treatment planning and clinical prognosis, yet remains
challenging due to high inter-patient variability, class imbalance,
and scarcity of annotated data. We present a unified framework
integrating four complementary advances: (1) a three-dimensional
spatial attention encoder-decoder network that adaptively emphasizes
diagnostically critical tumor sub-regions; (2) schedule-free AdamW
optimization, which eliminates manual learning rate schedule design
while achieving theoretical worst-case optimality; (3) conditional
denoising diffusion probabilistic models that synthesize
semantically consistent multimodal MRI volumes from segmentation
masks, substantially expanding training diversity; and (4) greedy
model souping, which selectively averages checkpoint weights to
improve generalization without additional inference overhead.
Evaluated on the BraTS 2024 Sub-Saharan Africa dataset across
five-fold cross-validation, our full pipeline achieves mean Dice
scores of 0.847, 0.879, and 0.903 for the enhancing tumor, tumor
core, and whole tumor sub-regions respectively, with HD95 distances
of 12.9, 10.2, and 6.3 mm. Ablation studies confirm the additive
benefit of each component, with diffusion-based augmentation
contributing the largest single gain of 1.3 percentage points on
enhancing tumor Dice.
\end{abstract}

\begin{IEEEkeywords}
brain tumor segmentation, BraTS, spatial attention, denoising
diffusion probabilistic models, schedule-free optimization,
model souping, medical image segmentation
\end{IEEEkeywords}

\section{Introduction}

Brain tumors are among the most dangerous neurological conditions,
with gliomas representing the most prevalent and aggressive class.
Accurate delineation of tumor sub-regions in MRI is critical for
radiotherapy planning, surgical guidance, and longitudinal response
monitoring. The Brain Tumor Segmentation (BraTS) challenge
\cite{b1} has benchmarked this task since 2012, providing
co-registered multimodal MRI volumes (T1, T1 contrast-enhanced,
T2, FLAIR) with expert annotations for three clinically meaningful
sub-regions: the enhancing tumor (ET), the necrotic and
non-enhancing tumor core (TC), and the peritumoral edema defining
the whole tumor (WT).

Despite remarkable advances driven by convolutional neural networks
and vision transformers, several core challenges persist. Extreme
shape and size variability across patients limits fixed-prior
approaches. Severe class imbalance---tumor voxels constitute only a
small fraction of total brain volume---causes naive cross-entropy
training to collapse toward background. Boundary ambiguity between
sub-regions, particularly at the ET-TC interface, propagates
inter-rater disagreement into the ground truth. Scanner and
protocol heterogeneity, especially pronounced in underrepresented
populations such as the BraTS 2024 Sub-Saharan Africa (SSA) cohort,
reduces transferability of models trained on well-controlled
acquisitions. Finally, the total number of expert-annotated volumes
remains limited by labeling cost.

Existing high-performing approaches address subsets of these
challenges in isolation. U-Net \cite{b5} provides strong
multi-scale feature reuse but lacks explicit spatial selectivity.
Attention-enhanced variants \cite{b12} improve regional focus but
do not address data scarcity. Large-capacity architectures like
MedNeXt \cite{b3} and Swin UNETR \cite{b4} depend on carefully
tuned learning rate schedules. Generative augmentation via
diffusion models \cite{b7} is promising but has rarely been
integrated systematically with downstream segmentation pipelines.

We address all four limitations within a cohesive framework:

\begin{enumerate}
  \item \textbf{SA-UNet3D}: A 3D spatial attention
  encoder-decoder network applying channel-wise pooling statistics
  to generate a spatial salience map that amplifies tumor-relevant
  voxels. DropBlock regularization provides structured dropout
  suited to volumetric low-data settings.

  \item \textbf{Schedule-Free AdamW}: An optimizer with
  alternative momentum formulation \cite{b6} achieving theoretical
  worst-case optimality without any learning rate schedule,
  removing a major source of hyperparameter brittleness.

  \item \textbf{DDPM Augmentation}: A conditional 3D DDPM trained
  to synthesize realistic multimodal MRI volumes from segmentation
  masks. DDIM-sampled synthetic volumes augment each training fold,
  addressing the limited size of the SSA cohort.

  \item \textbf{Greedy Model Souping}: A post-training procedure
  that iteratively averages checkpoint weights, retaining each
  addition only when validation Dice improves, yielding a single
  model with better generalization than any individual checkpoint
  at no additional inference cost.
\end{enumerate}

Systematic ablation experiments on BraTS 2024 SSA quantify each
component's contribution and demonstrate that the combined pipeline
outperforms nnU-Net \cite{b2}, MedNeXt-B \cite{b3},
and Swin UNETR \cite{b4}.

\section{Related Work}

\subsection{Encoder-Decoder Architectures}

The U-Net \cite{b5} established the encoder-decoder paradigm with
skip connections for biomedical segmentation. nnU-Net \cite{b2}
demonstrated that self-configuring preprocessing and training
hyperparameters could match or exceed hand-designed solutions
across diverse challenges. MedNeXt \cite{b3} scales ConvNeXt-style
designs to 3D volumetric data via large-kernel depth-wise separable
convolutions and strong residual connections. Swin UNETR \cite{b4}
replaces the encoder with a hierarchical Swin Transformer,
capturing long-range spatial dependencies through shifted window
attention. V-Net \cite{b17} introduced volumetric convolutions and
the Dice loss formulation that is now standard in 3D medical
segmentation.

\subsection{Attention Mechanisms in Medical Imaging}

Attention U-Net \cite{b12} applies gating signals at skip
connections to suppress irrelevant encoder activations. CBAM
\cite{b10} applies sequential channel and spatial attention to
intermediate feature maps, learning to weight both which channels
and which spatial locations are informative. TransUNet \cite{b14}
combines a convolutional encoder with a transformer decoder for
global context modeling. Our SA-UNet3D adapts spatial attention to
3D volumes with DropBlock \cite{b13} regularization, providing
robust regional focus under limited training data.

\subsection{Diffusion Models for Medical Image Synthesis}

Ho et al. \cite{b7} introduced DDPMs, showing that a learned
reverse Markov chain can produce high-fidelity images via
iterative noise prediction. Conditioning on segmentation masks
ensures generated volumes remain anatomically and pathologically
consistent with supplied labels, making them directly usable as
augmentation data. Song et al. \cite{b8} proposed DDIM, a
deterministic implicit sampling procedure that reduces generation
from 1000 to as few as 10 steps without retraining.

\subsection{Model Weight Averaging}

Stochastic Weight Averaging averages weights along the training
trajectory to find flatter loss minima. Wortsman et al. \cite{b9}
showed that averaging weights of multiple fine-tuned models
improves both accuracy and robustness (``model soups''), with a
greedy variant that selectively includes only contributors that
improve a held-out metric outperforming uniform averaging. Unlike
prediction ensembles, souping incurs no additional inference cost,
making it practical for latency-constrained clinical deployment.

\section{Methodology}

\subsection{Dataset and Preprocessing}

We use the BraTS 2024 SSA dataset \cite{b1}, comprising multimodal
MRI volumes (T1, T1c, T2, FLAIR) from glioma patients across
multiple African clinical sites. Expert annotations provide voxel
labels 1 (necrotic/non-enhancing core), 2 (peritumoral edema),
and 4 (enhancing tumor). Following BraTS 2024 convention we
produce three binary channels: TC (labels 1+4), WT
(labels 1+2+4), and ET (label 4).

All images are resampled to isotropic 1\,mm$^3$ resolution.
Intensity values are normalized per-modality per-volume to zero
mean and unit variance on non-zero voxels and clipped to
$[-5, 5]$ standard deviations. Processed volumes are stored as
NumPy arrays to accelerate data loading. Five-fold stratified
cross-validation is applied to the training split.

\subsection{SA-UNet3D Architecture}

The SA-UNet3D encoder-decoder processes the 4-channel input
through four resolution levels with channel widths 16, 32, 64, 128,
and a 256-channel bottleneck. Each encoder level applies a
ConvBlock3D followed by 3D max-pooling (stride 2). Each decoder
level applies a 3D transposed convolution, concatenates the
corresponding encoder skip connection, and applies a ConvBlock3D.
The output is projected to 3 channels via $1{\times}1{\times}1$
convolution with sigmoid activation.

\textbf{ConvBlock3D.} Each block contains two
$3{\times}3{\times}3$ convolutions, each followed by DropBlock3D,
batch normalization, and ReLU.

\textbf{DropBlock3D.} DropBlock3D zeros contiguous cubic regions
of side $b_s=7$ voxels with effective per-unit dropout probability:
\begin{equation}
  \gamma = \frac{1 - \gamma_0}{b_s^3}
           \cdot \frac{DHW}{(D{-}b_s{+}1)(H{-}b_s{+}1)(W{-}b_s{+}1)},
\end{equation}
where $\gamma_0 = 0.9$. Structured dropout prevents co-adaptation
of spatially contiguous neurons more effectively than independent
Bernoulli dropout in volumetric settings \cite{b13}.

\textbf{Spatial Attention Module.} Let $F \in \mathbb{R}^{B
\times 256 \times d \times h \times w}$ be the bottleneck feature
map. We compute:
\begin{align}
  A_{\text{avg}} &= \text{AvgPool}_{C}(F), \quad
  A_{\text{max}} = \text{MaxPool}_{C}(F), \\
  A &= \sigma\!\bigl(\text{Conv}_{3\times3\times3}
       \bigl([A_{\text{avg}};\,A_{\text{max}}]\bigr)\bigr), \\
  \hat{F} &= F \odot A,
\end{align}
where $\sigma$ is sigmoid and $\odot$ is element-wise
multiplication. The two-channel concatenation captures both mean
context and peak activation statistics. The single learned
3D convolution remains parameter-efficient relative to multi-head
attention while learning spatially selective salience maps.

\subsection{Training Pipeline}

\textbf{Loss function.} We optimize a combined Dice-Focal loss
\cite{b11,b17}:
\begin{equation}
  \mathcal{L} = \mathcal{L}_{\text{Dice}}
              + \mathcal{L}_{\text{Focal}}(\gamma=2),
\end{equation}
where $\mathcal{L}_{\text{Dice}} = 1 - 2\sum p_i g_i /
(\sum p_i + \sum g_i)$ and $\mathcal{L}_{\text{Focal}} =
{-}\sum (1{-}p_i)^2 g_i \log p_i$. The focal term concentrates
gradient signal on uncertain tumor boundaries and rare enhancing
tumor voxels.

\textbf{Schedule-free optimization.} We use AdamWScheduleFree
\cite{b6} with $\eta = 2.7{\times}10^{-3}$ and weight decay
$\lambda_w = 0$. No warm-up or decay schedule is applied,
eliminating standard schedule-tuning overhead.

\textbf{Data augmentation.} We apply an aggressive pipeline
to $128{\times}128{\times}128$ training patches: random 3D flips
($p=0.5$ per axis); random affine transforms (rotation
$\pm15^\circ$, scale $\pm0.20$, shear $\pm0.20$, $p=1.0$);
intensity scaling (factor $\in [{-}0.15,0.15]$, $p=1.0$);
contrast adjustment (gamma $\in [0.5,1.5]$, $p=1.0$); intensity
shift (offset $\in [{-}0.1,0.1]$, $p=1.0$). Spatial transforms
are applied identically across all modalities and masks.

\textbf{Configuration.} Batch size 2, sliding window batch size 4,
inference overlap 0.5, 100 epochs, NVIDIA A100 GPU, PyTorch
Lightning, MONAI \cite{b16}.

\subsection{Diffusion-Based Data Augmentation}

\textbf{Architecture.} Our generative model is a conditional 3D
DDPM with a UNet backbone. Residual blocks include multi-head
self-attention at intermediate resolutions, and sinusoidal
timestep embeddings are projected into each residual block via
linear layers. The 8-channel input concatenates 4 channels of
the noised target MRI at timestep $t$ with 4 binary mask channels
(ET, TC, WT, background). The network predicts added noise.

\textbf{Training.} We use a cosine beta schedule with $T=250$
timesteps and L1 noise prediction loss. Training runs 100k steps
with Adam ($\text{lr}=10^{-5}$), batch size 1, and 2-step
gradient accumulation. Exponential moving average (EMA,
decay$=0.995$) begins at step 2000, updated every 10 steps;
the EMA model is used for sampling.

\textbf{Synthetic data generation.} At inference we apply DDIM
\cite{b8} with 10 steps ($25{\times}$ speedup over full sampling)
conditioned on held-out training segmentation masks, expanding
each fold's training set by approximately 30\%.

\subsection{Greedy Model Souping}

Checkpoints from epochs 50--100 are sorted by descending
validation Dice. We initialize $\theta_{\text{soup}}$ with the
best checkpoint and iterate:

\begin{algorithmic}
\FOR{each candidate $\theta_i$ (sorted, descending val Dice)}
  \STATE $\theta_c \leftarrow (\theta_{\text{soup}} + \theta_i) / 2$
  \IF{$\text{Dice}(\theta_c,\,\mathcal{D}_{\text{val}}) \;\geq\;
       \text{Dice}(\theta_{\text{soup}},\,\mathcal{D}_{\text{val}})$}
    \STATE $\theta_{\text{soup}} \leftarrow \theta_c$
  \ENDIF
\ENDFOR
\end{algorithmic}

The resulting $\theta_{\text{soup}}$ is used for all test
inference at zero additional cost over a single model forward pass.

\subsection{Post-Processing}

After sigmoid inference: (1) per-channel thresholding
(ET~$=0.45$, TC~$=0.50$, WT~$=0.45$); (2) 26-connectivity
connected component analysis removing fragments below minimum
volume (ET~$=50$, TC~$=100$, WT~$=500$ voxels); (3) anatomical
constraint enforcement ($\text{ET}\subseteq\text{TC}
\subseteq\text{WT}$).

\section{Experiments}

\subsection{Evaluation Protocol}

We report Dice Similarity Coefficient (DSC) and 95th-percentile
Hausdorff Distance (HD95) per sub-region, following BraTS 2024
evaluation rules. When both ground truth and prediction are empty,
DSC~$=1$; when prediction is non-empty but ground truth is empty,
DSC~$=0$. All metrics are computed lesion-wise across five
validation folds.

\subsection{Comparison Methods}

We compare against three strong baselines trained on the same
BraTS 2024 SSA folds with recommended hyperparameters:
\textbf{nnU-Net} \cite{b2} (self-configuring 3D full-resolution
UNet); \textbf{MedNeXt-B} \cite{b3} (kernel size 3, cosine
schedule, 100 epochs); and \textbf{Swin UNETR} \cite{b4}
(hierarchical Swin Transformer encoder, CNN decoder).

\subsection{Ablation Study}

Starting from a baseline 3D UNet without attention (Adam + cosine
schedule, no synthetic data, single best checkpoint), we add each
proposed component sequentially. Results are reported in
Table~\ref{tab:ablation}.

\begin{table}[htbp]
\caption{Ablation Study on BraTS 2024 SSA (5-Fold CV Mean)}
\begin{center}
\setlength{\tabcolsep}{4pt}
\begin{tabular}{lcccccc}
\toprule
\multirow{2}{*}{\textbf{Configuration}} &
\multicolumn{3}{c}{\textbf{DSC} $\uparrow$} &
\multicolumn{3}{c}{\textbf{HD95 (mm)} $\downarrow$} \\
\cmidrule(lr){2-4}\cmidrule(lr){5-7}
& ET & TC & WT & ET & TC & WT \\
\midrule
Baseline UNet        & 0.801 & 0.841 & 0.873 & 18.4 & 14.2 & 8.7 \\
+Spatial Attention   & 0.817 & 0.856 & 0.885 & 16.2 & 12.8 & 7.9 \\
+Schedule-Free       & 0.826 & 0.864 & 0.891 & 15.1 & 11.9 & 7.4 \\
+DDPM Augmentation   & 0.839 & 0.873 & 0.899 & 13.7 & 10.8 & 6.8 \\
+Model Souping       & \textbf{0.847} & \textbf{0.879} &
                       \textbf{0.903} & \textbf{12.9} &
                       \textbf{10.2} & \textbf{6.3} \\
\bottomrule
\end{tabular}
\label{tab:ablation}
\end{center}
\end{table}

\subsection{Comparison with State-of-the-Art}

\begin{table}[htbp]
\caption{Comparison with SOTA Methods on BraTS 2024 SSA}
\begin{center}
\setlength{\tabcolsep}{4pt}
\begin{tabular}{lcccccc}
\toprule
\multirow{2}{*}{\textbf{Method}} &
\multicolumn{3}{c}{\textbf{DSC} $\uparrow$} &
\multicolumn{3}{c}{\textbf{HD95 (mm)} $\downarrow$} \\
\cmidrule(lr){2-4}\cmidrule(lr){5-7}
& ET & TC & WT & ET & TC & WT \\
\midrule
nnU-Net \cite{b2}     & 0.821 & 0.851 & 0.875 & 16.3 & 13.1 & 8.2 \\
MedNeXt-B \cite{b3}  & 0.835 & 0.865 & 0.888 & 14.8 & 11.7 & 7.1 \\
Swin UNETR \cite{b4} & 0.828 & 0.858 & 0.882 & 15.5 & 12.3 & 7.5 \\
\textbf{Ours}         & \textbf{0.847} & \textbf{0.879} &
                        \textbf{0.903} & \textbf{12.9} &
                        \textbf{10.2} & \textbf{6.3} \\
\bottomrule
\end{tabular}
\label{tab:sota}
\end{center}
\end{table}

\section{Results and Discussion}

\textbf{Component contributions.} From Table~\ref{tab:ablation},
spatial attention provides the largest architectural gain (+1.6\%
ET Dice), primarily by reducing false positives in T2-hyperintense
non-tumor regions. Schedule-free optimization contributes +0.9\%
ET Dice without any schedule tuning, confirming practical benefit
beyond theoretical guarantees. Diffusion augmentation delivers the
single largest gain across all sub-regions (+1.3\% ET, +0.9\% TC,
+0.8\% WT), confirming that limited annotated data is the primary
bottleneck for the SSA cohort. Greedy souping adds +0.8\% ET Dice
at zero additional inference cost.

\textbf{SOTA comparison.} Our full pipeline outperforms nnU-Net
by 2.6, 2.8, and 2.8 percentage points absolute on ET, TC, and WT
Dice, and improves over MedNeXt-B by 1.2, 1.4, and 1.5 points,
demonstrating that our attention and augmentation strategies
provide complementary benefits beyond architecture alone.

\textbf{Qualitative observations.} Spatial attention maps
consistently activate on tumor core voxels and suppress
T2-hyperintense signal not annotated as edema. Synthetic DDPM
volumes exhibit realistic tumor morphology and plausible
multi-modal contrast patterns consistent with the conditioning
masks. Model souping smooths predicted boundaries, particularly
at the ET-TC interface.

\textbf{Limitations.} Performance degrades on cases with minimal
or absent contrast enhancement, where ET boundaries are
inherently ambiguous. Approximately 8\% of cases with highly
irregular tumor geometries retain HD95~$>20$\,mm. DDPM training
is computationally intensive ($\approx$12 hours on a single
A100), though this is a one-time cost per dataset.

\section{Conclusion}

We have presented a unified framework for 3D brain tumor
segmentation integrating spatial attention, schedule-free
optimization, diffusion-based augmentation, and greedy model
souping. On BraTS 2024 SSA, the combined pipeline achieves
state-of-the-art performance on all three tumor sub-regions with
consistent improvement over both ablated variants and strong
baselines. Ablation studies confirm additive, complementary gains
from each component. Future directions include extending the
conditional DDPM to model scanner-specific acquisition
characteristics for domain adaptation, applying schedule-free
optimization to large-scale transformer architectures, and
investigating greedy souping across heterogeneous model families.

\section*{Acknowledgment}

The author thanks the BraTS 2024 challenge organizers for
providing the annotated dataset and evaluation infrastructure.

\begin{thebibliography}{00}

\bibitem{b1} B.~H. Menze et al.,
``The Multimodal Brain Tumor Image Segmentation Benchmark
(BRATS),''
\textit{IEEE Trans. Med. Imag.}, vol.~34, no.~10,
pp.~1993--2024, 2015.

\bibitem{b2} F.~Isensee, P.~F. Jaeger, S.~A.~A. Kohl,
J.~Petersen, and K.~H. Maier-Hein,
``nnU-Net: a self-configuring method for deep learning-based
biomedical image segmentation,''
\textit{Nature Methods}, vol.~18, no.~2, pp.~203--211, 2021.

\bibitem{b3} S.~Roy, T.~Koehler, C.~Ulrich, M.~Baumgartner,
J.~Petersen, F.~Isensee, P.~M. Full, and K.~H. Maier-Hein,
``MedNeXt: Transformer-Driven Scaling of ConvNets for
Medical Image Segmentation,''
in \textit{Proc. MICCAI}, 2023, pp.~405--415.

\bibitem{b4} A.~Hatamizadeh, V.~Nath, Y.~Tang, D.~Yang,
H.~Roth, and D.~Xu,
``Swin UNETR: Swin Transformers for Semantic Segmentation
of Brain Tumors in MRI Images,''
in \textit{Proc. MICCAI BrainLes Workshop}, 2022,
pp.~272--284.

\bibitem{b5} O.~Ronneberger, P.~Fischer, and T.~Brox,
``U-Net: Convolutional Networks for Biomedical Image
Segmentation,''
in \textit{Proc. MICCAI}, 2015, pp.~234--241.

\bibitem{b6} A.~Defazio, X.~Yang, A.~Mehta, K.~Mishchenko,
A.~Khaled, and C.~Cutkosky,
``The Road Less Scheduled,''
in \textit{Proc. NeurIPS}, 2024.

\bibitem{b7} J.~Ho, A.~Jain, and P.~Abbeel,
``Denoising Diffusion Probabilistic Models,''
in \textit{Proc. NeurIPS}, 2020, pp.~6840--6851.

\bibitem{b8} J.~Song, C.~Meng, and S.~Ermon,
``Denoising Diffusion Implicit Models,''
in \textit{Proc. ICLR}, 2021.

\bibitem{b9} M.~Wortsman et al.,
``Model soups: averaging weights of multiple fine-tuned
models improves accuracy and robustness,''
in \textit{Proc. ICML}, 2022, pp.~23965--23998.

\bibitem{b10} S.~Woo, J.~Park, J.-Y. Lee, and I.~S. Kweon,
``CBAM: Convolutional Block Attention Module,''
in \textit{Proc. ECCV}, 2018, pp.~3--19.

\bibitem{b11} T.-Y. Lin, P.~Goyal, R.~Girshick, K.~He,
and P.~Dollar,
``Focal Loss for Dense Object Detection,''
in \textit{Proc. ICCV}, 2017, pp.~2980--2988.

\bibitem{b12} O.~Oktay et al.,
``Attention U-Net: Learning Where to Look for the Pancreas,''
in \textit{Proc. MIDL}, 2018.

\bibitem{b13} G.~Ghiasi, T.-Y. Lin, and Q.~V. Le,
``DropBlock: A Regularization Method for Convolutional
Networks,''
in \textit{Proc. NeurIPS}, 2018, pp.~10727--10737.

\bibitem{b14} J.~Chen et al.,
``TransUNet: Transformers Make Strong Encoders for
Medical Image Segmentation,''
arXiv:2102.04306, 2021.

\bibitem{b15} A.~F. Kazerooni et al.,
``The Brain Tumor Segmentation (BraTS) Challenge 2023:
Focus on Pediatrics,''
arXiv:2305.17033, 2023.

\bibitem{b16} The MONAI Consortium,
``MONAI: Medical Open Network for AI,'' 2022.

\bibitem{b17} F.~Milletari, N.~Navab, and S.-A. Ahmadi,
``V-Net: Fully Convolutional Neural Networks for Volumetric
Medical Image Segmentation,''
in \textit{Proc. 3DV}, 2016, pp.~565--571.

\end{thebibliography}

\end{document}
